\chapter{2015年重庆卷（理）}

\section{单选题}

\begin{problem}
已知集合 \(A = \{1, 2, 3\}\)，\(B = \{2, 3\}\)，则（\quad）

\choices
    {\(A = B\)}
    {\(A \cap B = \varnothing\)}
    {\(A \not\subset B\)}
    {\(B \subsetneq A\)}

\begin{answer}
D
\end{answer}

\begin{solution}
因为 \(1\in A\) 而 \(1\notin B\)，所以 \(A\ne B\)，且 \(A\not\subset B\). 又因为 \(A\cap B=B=\{2,3\}\ne\varnothing\)，选项 B 错误；又因 \(B\) 中的每个元素都属于 \(A\)，并且 \(1\in A\setminus B\)，故 \(B\subsetneq A\)，所以选项 D 正确.
\end{solution}
\end{problem}

\begin{problem}
在等差数列 \(\{a_{n}\}\) 中，若 \(a_{2} = 4\)，\(a_{4} = 2\)，则 \(a_{6} =\)（\quad）

\choices
    {\(-1\)}
    {\(0\)}
    {\(1\)}
    {\(6\)}

\begin{answer}
B
\end{answer}

\begin{solution}
设等差数列的公差为 \(d\)，则 \(a_{4}-a_{2}=2d\)，所以
\[
d=\frac{2-4}{2}=-1
\]
于是 \(a_{6}=a_{4}+2d=2-2=0\)，故答案为 B.
\end{solution}
\end{problem}

\begin{problem}
重庆市 2013 年各月的平均气温（\(^{\circ}\mathrm{C}\)） 数据的茎叶图如下：

\begin{center}
\renewcommand{\arraystretch}{1.2}
\begin{tabular}{r|l}
0 & 8\quad 9 \\
1 & 2\quad 5\quad 8 \\
2 & 0\quad 0\quad 3\quad 3\quad 8 \\
3 & 1\quad 2
\end{tabular}
\end{center}

则这组数据的中位数是（\quad）

\choices
    {\(19\)}
    {\(20\)}
    {\(21.5\)}
    {\(23\)}

\begin{answer}
B
\end{answer}

\begin{solution}
由茎叶图，这组数据按从小到大排列为 \(8\)，\(9\)，\(12\)，\(15\)，\(18\)，\(20\)，\(20\)，\(23\)，\(23\)，\(28\)，\(31\)，\(32\)，共有 \(12\) 个数. 因此中位数为第 \(6\) 个数和第 \(7\) 个数的平均数，即
\[
\frac{20+20}{2}=20
\]
故答案为 B.
\end{solution}
\end{problem}

\begin{problem}
“\(x > 1\)”是“\(\log_{\frac{1}{2}}(x + 2) < 0\)”的（\quad）

\choices
    {充要条件}
    {充分不必要条件}
    {必要而不充分条件}
    {既不充分也不必要条件}

\begin{answer}
B
\end{answer}

\begin{solution}
因为对数函数 \(y=\log_{\frac12}x\) 在 \((0,+\infty)\) 上为减函数，且其真数须为正，所以
\[
\log_{\frac12}(x+2)<0\iff x+2>1\iff x>-1
\]
由 \(x>1\) 一定能推出 \(x>-1\)，但 \(x=0\) 满足 \(x>-1\) 而不满足 \(x>1\). 因此“\(x>1\)”是题中不等式的充分不必要条件，故答案为 B.
\end{solution}
\end{problem}

\begin{problem}
某几何体的三视图如图所示，则该几何体的体积为（\quad）

\bitmapfigure[max width=0.42\linewidth]{img/2015/chongqing_science/q5_fig1.png}

\choices
    {\(\frac{1}{3} + \pi\)}
    {\(\frac{2}{3} + \pi\)}
    {\(\frac{1}{3} + 2\pi\)}
    {\(\frac{2}{3} + 2\pi\)}

\begin{answer}
A
\end{answer}

\begin{solution}
由左视图可知，右侧部分是底面半径为 \(1\)、高为 \(2\) 的半圆柱；左侧部分是高为 \(1\) 的三棱锥，其底面三角形的底为 \(2\)、高为 \(1\)，底面积为 \(1\). 因此
\[
V=\frac13\left(\frac12\times1\times2\right)\times1+\frac12\pi\times1^2\times2=\frac13+\pi
\]
故答案为 A.
\end{solution}
\end{problem}

\begin{problem}
若非零向量 \(\bs{a}\)，\(\bs{b}\) 满足 \( \left|\bs{a}\right|=\frac{2\sqrt{2}}{3}\left|\bs{b}\right| \)，且 \( (\bs{a}-\bs{b})\perp(3\bs{a}+2\bs{b}) \)，则 \(\bs{a}\) 与 \(\bs{b}\) 的夹角为（\quad）

\choices
    {\( \frac{\pi}{4} \)}
    {\( \frac{\pi}{2} \)}
    {\( \frac{3\pi}{4} \)}
    {\(\pi\)}

\begin{answer}
A
\end{answer}

\begin{solution}
设 \(\bs{a}\) 与 \(\bs{b}\) 的夹角为 \(\theta\)，并记 \(\left|\bs{b}\right|=m>0\). 由已知
\[
\left|\bs{a}\right|^2=\frac89m^2
\]
又由垂直条件，
\[
0=(\bs{a}-\bs{b})\cdot(3\bs{a}+2\bs{b})=3\left|\bs{a}\right|^2-\bs{a}\cdot\bs{b}-2\left|\bs{b}\right|^2
\]
所以
\[
\bs{a}\cdot\bs{b}=3\times\frac89m^2-2m^2=\frac23m^2
\]
于是
\[
\cos\theta=\frac{\bs{a}\cdot\bs{b}}{\left|\bs{a}\right|\left|\bs{b}\right|}=\frac{\frac23m^2}{\frac{2\sqrt2}{3}m^2}=\frac{\sqrt2}{2}
\]
因为 \(0\leqslant\theta\leqslant\pi\)，故 \(\theta=\frac\pi4\)，答案为 A.
\end{solution}
\end{problem}

\begin{problem}
执行如图所示的程序框图，若输出 \(k\) 的值为 \(8\)，则判断框内可填入的条件是（\quad）

\bitmapfigure[max width=0.50\linewidth]{img/2015/chongqing_science/q7_fig1.png}

\choices
    {\(s \leqslant \frac{3}{4}\)}
    {\(s \leqslant \frac{5}{6}\)}
    {\(s \leqslant \frac{11}{12}\)}
    {\(s \leqslant \frac{25}{24}\)}

\begin{answer}
C
\end{answer}

\begin{solution}
程序每次先将 \(k\) 增加 \(2\)，再把 \(\frac1k\) 加到 \(s\) 中. 若最终输出 \(k=8\)，则循环应在 \(k=2\)，\(4\)，\(6\) 时继续，在 \(k=8\) 时停止. 相应地，
\(s_2=\frac12\)，\(s_4=\frac12+\frac14=\frac34\)，\(s_6=\frac34+\frac16=\frac{11}{12}\)
而再执行一次得到
\[
s_8=\frac{11}{12}+\frac18=\frac{25}{24}>\frac{11}{12}
\]
因此判断条件应为 \(s\leqslant\frac{11}{12}\)，故答案为 C.
\end{solution}
\end{problem}

\begin{problem}
已知直线 \(l: x + ay - 1 = 0 \) （\(a \in \R\)）是圆 \(C: x^2 + y^2 - 4x - 2y + 1 = 0\) 的对称轴. 过点 \(A(-4, a)\) 作圆 \(C\) 的一条切线. 切点为 \(B\)，则 \(|AB| =\)（\quad）

\choices
    {\(2\)}
    {\(4\sqrt{2}\)}
    {\(6\)}
    {\(2\sqrt{10}\)}

\begin{answer}
C
\end{answer}

\begin{solution}
将圆方程配方得
\[
(x-2)^2+(y-1)^2=4
\]
所以圆心为 \(O(2,1)\)，半径为 \(2\). 圆的对称轴必须经过圆心，将 \(O\) 代入直线方程得 \(2+a-1=0\)，故 \(a=-1\). 于是 \(A(-4,-1)\)，从而
\[
OA=\sqrt{(-4-2)^2+(-1-1)^2}=2\sqrt{10}
\]
在切点 \(B\) 处，\(OB\perp AB\)，所以
\[
AB=\sqrt{OA^2-OB^2}=\sqrt{40-4}=6
\]
故答案为 C.
\end{solution}
\end{problem}

\begin{problem}
若 \(\tan \alpha = 2\tan \frac{\pi}{5}\)，则 \(\frac{\cos\left(\alpha - \frac{3\pi}{10}\right)}{\sin\left(\alpha - \frac{\pi}{5}\right)} =\)（\quad）

\choices
    {\(1\)}
    {\(2\)}
    {\(3\)}
    {\(4\)}

\begin{answer}
C
\end{answer}

\begin{solution}
记 \(\beta=\frac\pi5\)，则 \(\frac{3\pi}{10}=\frac\pi2-\beta\). 利用 \(\tan\alpha=2\tan\beta\) 及两角和差公式，得
\[
\cos\left(\alpha-\frac{3\pi}{10}\right)=\cos\alpha\sin\beta+\sin\alpha\cos\beta
\]
\[
\sin\left(\alpha-\beta\right)=\sin\alpha\cos\beta-\cos\alpha\sin\beta
\]
又因为 \(\sin\alpha=2\tan\beta\cos\alpha\)，所以
\(\cos\left(\alpha-\frac{3\pi}{10}\right)=3\sin\beta\cos\alpha\)，\(\sin\left(\alpha-\beta\right)=\sin\beta\cos\alpha\).
分母不为零，故所求比值为 \(3\)，答案为 C.
\end{solution}
\end{problem}

\begin{problem}
设双曲线 \(\frac{x^2}{a^2} - \frac{y^2}{b^2} = 1 \) （\(a > 0\)，\(b > 0\)）的右焦点为 \(F\)，右顶点为 \(A\)，过 \(F\) 作 \(AF\) 的垂线与双曲线交于 \(B\)，\(C\) 两点，过 \(B\)，\(C\) 分别作 \(AC\)，\(AB\) 的垂线，两垂线交于点 \(D\). 若 \(D\) 到直线 \(BC\) 的距离小于 \(a + \sqrt{a^2 + b^2}\)，则该双曲线的渐近线斜率的取值范围是（\quad）

\choices
    {\((-1,0)\cup (0,1)\)}
    {\((-\infty, -1) \cup (1, +\infty)\)}
    {\((-\sqrt{2}, 0) \cup (0, \sqrt{2})\)}
    {\((-\infty, -\sqrt{2}) \cup (\sqrt{2}, +\infty)\)}

\begin{answer}
A
\end{answer}

\begin{solution}
记 \(c=\sqrt{a^2+b^2}\)，则右焦点和右顶点分别为 \(F(c,0)\)、\(A(a,0)\). 直线 \(BC\) 为 \(x=c\)，由双曲线方程得交点的纵坐标满足
\(\frac{c^2}{a^2}-\frac{y^2}{b^2}=1\)，\(\Longrightarrow\)，\(y^2=\frac{b^4}{a^2}\).
取 \(q=\frac{b^2}{a}>0\)，于是
\(B(c,q)\)，\(C(c,-q)\).
直线 \(AC\) 的斜率为 \(-\frac{q}{c-a}\)，直线 \(AB\) 的斜率为 \(\frac{q}{c-a}\)，其中 \(c-a>0\). 因此过 \(B\) 作 \(AC\) 的垂线、过 \(C\) 作 \(AB\) 的垂线，交点 \(D\) 满足
\(y-q=\frac{c-a}{q}(x-c)\)，\(y+q=-\frac{c-a}{q}(x-c)\).
两式相加得 \(y=0\)，再代入任一式得 \(x-c=-\frac{q^2}{c-a}\)，所以
\[
d(D,BC)=\frac{q^2}{c-a}=\frac{b^4/a^2}{c-a}=\frac{b^2(a+c)}{a^2}
\]
令 \(t=\frac ca>1\)，则 \(\frac{b^2}{a^2}=t^2-1\)，题设不等式化为
\[
a(t^2-1)(t+1)<a(t+1)
\]
即 \(1<t<\sqrt2\). 双曲线渐近线斜率为 \(\pm\frac ba=\pm\sqrt{t^2-1}\)，故其取值范围为 \((-1,0)\cup(0,1)\)，答案为 A.
\end{solution}
\end{problem}

\section{填空题}

\begin{problem}
设复数 \(a + b\i\) （\(a\)，\(b \in \R\)）的模为 \(\sqrt{3}\)，则 \((a + b\i)(a - b\i) = \)\fillinblank{}.

\begin{answer}
\(3\)
\end{answer}

\begin{solution}
复数的模满足 \(\left|a+b\i\right|^2=(a+b\i)(a-b\i)\). 已知 \(\left|a+b\i\right|=\sqrt3\)，所以
\[
(a+b\i)(a-b\i)=\left|a+b\i\right|^2=3
\]
\end{solution}
\end{problem}

\begin{problem}
\(\left(x^{3} + \frac{1}{2\sqrt{x}}\right)^{5}\) 的展开式中 \(x^{8}\) 的系数是\fillinblank{}. （用数字作答）

\begin{answer}
\(\frac{5}{2}\)
\end{answer}

\begin{solution}
设通项中取 \(\frac{1}{2\sqrt{x}}\) 的因子数为 \(k\)，则该项为
\[
\mathrm C_5^k\left(x^3\right)^{5-k}\left(\frac{1}{2\sqrt{x}}\right)^k
\]
其中 \(x\) 的指数为 \(3(5-k)-\frac{k}{2}=15-\frac{7k}{2}\). 令其等于 \(8\)，得 \(k=2\). 因此所求系数为
\[
\mathrm C_5^2\left(\frac12\right)^2=\frac{5}{2}
\]
\end{solution}
\end{problem}

\begin{problem}
在 \(\triangle ABC\) 中，\(B = 120^{\circ}\)，\(AB = \sqrt{2}\)，\(A\) 的角平分线 \(AD = \sqrt{3}\)，则 \(AC = \)\fillinblank{}.

\begin{answer}
\(\sqrt{6}\)
\end{answer}

\begin{solution}
设 \(AC=\sqrt2u\)，\(BC=\sqrt2v\)，其中 \(u\)，\(v>0\). 由 \(\angle B=120^\circ\) 及余弦定理，
\[
u^2=v^2+v+1
\]
由角平分线长公式
\[
AD^2=AB\cdot AC\left(1-\frac{BC^2}{(AB+AC)^2}\right)
\]
得
\[
3=2u\left(1-\frac{v^2}{(u+1)^2}\right)
\]
利用 \(v^2=u^2-v-1\) 化简上式，得
\(2uv=-u^2+2u+3\)，\(qquad v=\frac{-u^2+2u+3}{2u}\).
将此式代入 \(u^2=v^2+v+1\)，并通分化简得
\[
0=-\frac{3(u^2-3)(u+1)^2}{4u^2}
\]
由于 \(u>0\)，且 \(u+1\ne0\)，所以 \(u^2=3\)，从而 \(u=\sqrt3\). 此时由前式得 \(v=1>0\)，满足三角形条件，故 \(AC=\sqrt2u=\sqrt6\).
\end{solution}
\end{problem}

\section{选做题}

\begin{problem}
如图，圆 \(O\) 的弦 \(AB\)，\(CD\) 相交于点 \(E\)，过点 \(A\) 作圆 \(O\) 的切线与 \(DC\) 的延长线交于点 \(P\)，若 \(PA = 6\)，\(AE = 9\)，\(PC = 3\)，\(CE:ED = 2:1\)，则 \(BE = \)\fillinblank{}.

\bitmapfigure[max width=0.32\linewidth]{img/2015/chongqing_science/q14_fig1.png}

\begin{answer}
\(2\)
\end{answer}

\begin{solution}
由切线与割线定理，\(PA^2=PC\cdot PD\)，所以
\[
PD=\frac{PA^2}{PC}=\frac{36}{3}=12
\]
图中点的顺序为 \(P\)，\(C\)，\(E\)，\(D\)，故 \(CD=PD-PC=9\). 由 \(CE:ED=2:1\)，得 \(CE=6\)，\(ED=3\). 再由相交弦定理，
\[
AE\cdot BE=CE\cdot ED
\]
从而
\[
BE=\frac{6\times3}{9}=2
\]
\end{solution}
\end{problem}

\begin{problem}
已知直线 \(l\) 的参数方程为 \(\begin{cases}
x = -1 + t,\\
y = 1 + t,
\end{cases}\) （\(t\) 为参数），以坐标原点为极点，\(x\) 轴的正半轴为极轴建立极坐标系，曲线 \(C\) 的极坐标方程为 \(\rho^2 \cos 2\theta = 4\) （\(\rho > 0\)，\(\frac{3\pi}{4} < \theta < \frac{5\pi}{4}\)），则直线 \(l\) 与曲线 \(C\) 的交点的极坐标为\fillinblank{}.

\begin{answer}
\((2,\pi)\)
\end{answer}

\begin{solution}
由直线的参数方程，消去参数得 \(y-x=2\)，即
\[
\rho(\sin\theta-\cos\theta)=2
\]
又因为 \(\rho^2\cos2\theta=4\)，所以
\[
-\rho^2(\sin\theta-\cos\theta)(\sin\theta+\cos\theta)=4
\]
代入直线关系，得 \(\rho(\sin\theta+\cos\theta)=-2\). 两式相加得 \(2\rho\sin\theta=0\)，由于 \(\rho>0\) 且 \(\frac{3\pi}{4}<\theta<\frac{5\pi}{4}\)，故 \(\theta=\pi\)，再由直线关系得 \(\rho=2\). 所以交点的极坐标为 \((2,\pi)\).
\end{solution}
\end{problem}

\begin{problem}
若函数 \( f(x) = |x + 1| + 2|x - a| \) 的最小值为 5，则实数 \( a = \)\fillinblank{}.

\begin{answer}
\(-6\) 或 \(4\)
\end{answer}

\begin{solution}
当 \(a\geqslant-1\) 时，分三段展开：
\[
f(x)=
\begin{cases}
2a-3x-1, & x\leqslant-1,\\
2a+1-x, & -1\leqslant x\leqslant a,\\
3x+1-2a, & x\geqslant a.
\end{cases}
\]
前两段的斜率分别为 \(-3\)、\(-1\)，后一段的斜率为 \(3\)，所以函数先递减后递增，最小值在 \(x=a\) 处取得，且 \(f(a)=a+1\). 令 \(a+1=5\)，得 \(a=4\).

当 \(a\leqslant-1\) 时，分三段展开：
\[
f(x)=
\begin{cases}
2a-3x-1, & x\leqslant a,\\
x-1-2a, & a\leqslant x\leqslant-1,\\
3x+1-2a, & x\geqslant-1.
\end{cases}
\]
三段的斜率分别为 \(-3\)、\(1\)、\(3\)，所以函数先递减后递增，最小值在 \(x=a\) 处取得，且 \(f(a)=-a-1\). 令 \(-a-1=5\)，得 \(a=-6\).

因此 \(a=-6\) 或 \(a=4\).
\end{solution}
\end{problem}

\section{解答题}

\begin{problem}
端午节吃粽子是我国的传统习俗. 设一盘中装有 \(10\) 个粽子，其中豆沙粽 \(2\) 个，肉粽 \(3\) 个，白粽 \(5\) 个，这三种粽子的外观完全相同，从中任意选取 \(3\) 个.

\begin{enumerate}
\item 求三种粽子各取到 1 个的概率；
\item 设 \(X\) 表示取到的豆沙粽个数，求 \(X\) 的分布列与数学期望.
\end{enumerate}
\end{problem}

\begin{solution}
\begin{enumerate}
\item 从 \(10\) 个粽子中任取 \(3\) 个的基本事件总数为 \(\mathrm C_{10}^{3}\). 三种粽子各取到 \(1\) 个的事件数为 \(\mathrm C_2^1\mathrm C_3^1\mathrm C_5^1\)，所以所求概率为
\[
P=\frac{\mathrm C_2^1\mathrm C_3^1\mathrm C_5^1}{\mathrm C_{10}^{3}}=\frac{2\times3\times5}{120}=\frac14
\]
\item 因为豆沙粽只有 \(2\) 个，所以 \(X\) 的可能取值为 \(0\)，\(1\)，\(2\). 分别计算得
\[
P(X=0)=\frac{\mathrm C_8^3}{\mathrm C_{10}^{3}}=\frac{56}{120}=\frac{7}{15}
\]
\[
P(X=1)=\frac{\mathrm C_2^1\mathrm C_8^2}{\mathrm C_{10}^{3}}=\frac{56}{120}=\frac{7}{15}
\]
\[
P(X=2)=\frac{\mathrm C_2^2\mathrm C_8^1}{\mathrm C_{10}^{3}}=\frac{8}{120}=\frac{1}{15}
\]
因此 \(X\) 的分布列为
\begin{center}
\begin{tabular}{c|ccc}
\(X\) & \(0\) & \(1\) & \(2\) \\
\hline
\(P\) & \(\dfrac{7}{15}\) & \(\dfrac{7}{15}\) & \(\dfrac{1}{15}\)
\end{tabular}
\end{center}
且
\[
E(X)=0\times\frac{7}{15}+1\times\frac{7}{15}+2\times\frac{1}{15}=\frac35
\]
\end{enumerate}
\end{solution}

\begin{problem}
已知函数 \( f(x) = \sin \left( \frac{\pi}{2} - x \right) \sin x - \sqrt{3} \cos^2 x \).

\begin{enumerate}
\item 求 \( f(x) \) 的最小正周期和最大值；
\item 讨论 \( f(x) \) 在 \( \left[\frac{\pi}{6}, \frac{2\pi}{3}\right] \) 上的单调性.
\end{enumerate}
\end{problem}

\begin{solution}
\begin{enumerate}
\item 由诱导公式、二倍角公式，
\[
\begin{split}
f(x)&=\sin x\cos x-\sqrt3\cos^2x\\
&=\frac12\sin2x-\frac{\sqrt3}{2}(1+\cos2x)\\
&=\sin\left(2x-\frac\pi3\right)-\frac{\sqrt3}{2}.
\end{split}
\]
函数 \(\sin(2x-\frac\pi3)\) 的周期为 \(\pi\)，且不存在更小的正周期，因此 \(f(x)\) 的最小正周期为 \(\pi\). 因为正弦函数的最大值为 \(1\)，所以 \(f(x)\) 的最大值为 \(1-\frac{\sqrt3}{2}\).
\item 当 \(x\in\left[\frac\pi6,\frac{2\pi}{3}\right]\) 时，令 \(u=2x-\frac\pi3\)，则 \(u\in[0,\pi]\). 正弦函数在 \([0,\frac\pi2]\) 上单调递增，在 \([\frac\pi2,\pi]\) 上单调递减，而 \(u=\frac\pi2\) 对应 \(x=\frac{5\pi}{12}\). 因此 \(f(x)\) 在 \(\left[\frac\pi6,\frac{5\pi}{12}\right]\) 上单调递增，在 \(\left[\frac{5\pi}{12},\frac{2\pi}{3}\right]\) 上单调递减.
\end{enumerate}
\end{solution}

\begin{problem}
如图，三棱锥 \(P - ABC\) 中，\(PC \perp\) 平面 \(ABC\)，\(PC = 3\)，\(\angle ACB = \frac{\pi}{2}\)，\(D\)，\(E\) 分别为线段 \(AB\)，\(BC\) 上的点，且 \(CD = DE = \sqrt{2}\)，\(CE = 2EB = 2\).

\begin{enumerate}
\item 证明： \(DE \perp\) 平面 \(PCD\)；
\item 求二面角 \(A - PD - C\) 的余弦值.
\end{enumerate}

\FigureLayoutDeclare{img/2015/chongqing_science/q19_fig1.png}{9.5cm}{94bp 230bp 122bp 7bp}
\bitmapfigure[max width=0.34\linewidth]{img/2015/chongqing_science/q19_fig1.png}
\end{problem}

\begin{solution}
\begin{enumerate}
\item 因为 \(PC\perp\) 平面 \(ABC\)，且 \(DE\) 在平面 \(ABC\) 内，所以 \(PC\perp DE\). 又在三角形 \(CDE\) 中，
\[
CD^2+DE^2=2+2=4=CE^2
\]
故 \(CD\perp DE\). 直线 \(PC\)、\(CD\) 是平面 \(PCD\) 内相交的两条直线，因而 \(DE\perp\) 平面 \(PCD\).
\item 取 \(C\) 为原点，分别以 \(CB\)、\(CA\)、\(CP\) 所在直线为 \(x\)、\(y\)、\(z\) 轴的正方向. 由 \(CE=2\)、\(EB=1\)，得
\(E(2,0,0)\)，\(B(3,0,0)\)，\(P(0,0,3)\).
设 \(A(0,h,0)\)，因 \(D\in AB\)，可设 \(D=(3t,h(1-t),0)\). 由 \(CD=DE\)，
\[
9t^2+h^2(1-t)^2=(3t-2)^2+h^2(1-t)^2
\]
所以 \(t=\frac13\). 再由 \(CD^2=2\)，得 \(1+\frac49h^2=2\)，即 \(h=\frac32\). 故
\(D(1,1,0)\)，\(A\left(0,\frac32,0\right)\).
平面 \(APD\) 的一个法向量可取
\[
\overrightarrow{DA}\times\overrightarrow{DP}=\frac32(1,2,1)
\]
平面 \(CPD\) 的一个法向量可取 \(\overrightarrow{DP}\times\overrightarrow{DC}=3(1,-1,0)\). 因此二面角的余弦值为
\[
\frac{\left|(1,2,1)\cdot(1,-1,0)\right|}{\sqrt{1^2+2^2+1^2}\sqrt{1^2+(-1)^2}}=\frac{1}{\sqrt6\sqrt2}=\frac{\sqrt3}{6}
\]
\end{enumerate}
\end{solution}

\begin{problem}
设函数 \( f(x) = \frac{3x^2 + ax}{\e^x}  \)（\(a \in \R\)）.

\begin{enumerate}
\item 若 \( f(x) \) 在 \( x = 0 \) 处取得极值，确定 \( a \) 的值，并求此时曲线 \( y = f(x) \) 在点 \((1, f(1))\) 处的切线方程；
\item 若 \( f(x) \) 在 \([3, +\infty)\) 上为减函数，求 \( a \) 的取值范围.
\end{enumerate}
\end{problem}

\begin{solution}
\begin{enumerate}
\item 对函数求导得
\[
f'(x)=\frac{-3x^2+(6-a)x+a}{\e^x}
\]
函数在 \(x=0\) 处取得极值，必须有 \(f'(0)=0\)，故 \(a=0\). 此时 \(f(x)=3x^2\e^{-x}\)，且 \(f(x)\geqslant0=f(0)\)，所以确实在 \(x=0\) 处取得极小值. 此时
\(f(1)=\frac3\e\)，\(qquad f'(1)=\frac3\e\).
由点斜式，切线方程为 \(y-\frac3\e=\frac3\e(x-1)\)，即 \(3x-\e y=0\).
\item 因为 \(\e^x>0\)，所以 \(f(x)\) 在 \([3,+\infty)\) 上为减函数，当且仅当
\(q(x)=-3x^2+(6-a)x+a\leqslant0\)，\((x\geqslant3)\).
若 \(a\geqslant-12\)，则 \(q'(x)=6-a-6x\leqslant0\)（\(x\geqslant3\)），故只需且必须有
\[
q(3)=-9-2a\leqslant0
\]
即 \(a\geqslant-\frac92\). 若 \(a<-12\)，则二次函数 \(q\) 的顶点横坐标 \(1-\frac a6>3\)，且顶点函数值为
\[
q\left(1-\frac a6\right)=a+\frac{(6-a)^2}{12}=\frac{a^2}{12}+3>0
\]
不符合减函数条件. 因此 \(a\in\left[-\frac92,+\infty\right)\).
\end{enumerate}
\end{solution}

\begin{problem}
如图，椭圆 \(\frac{x^2}{a^2} + \frac{y^2}{b^2} = 1 \) （\(a > b > 0\)）的左，右焦点分别为 \(F_1\)，\(F_2\)，过 \(F_2\) 的直线交椭圆于 \(P\)，\(Q\) 两点，且 \(PQ \perp PF_1\).

\begin{enumerate}
\item 若 \( |PF_1| = 2 + \sqrt{2} \)，\( |PF_2| = 2 - \sqrt{2} \)，求椭圆标准方程；
\item 若 \( |PF_1| = |PQ| \)，求椭圆的离心率 \( e \).
\end{enumerate}

\bitmapfigure[max width=0.34\linewidth]{img/2015/chongqing_science/q21_fig1.png}
\end{problem}

\begin{solution}
\begin{enumerate}
\item 由 \(PQ\perp PF_1\) 且 \(P\)，\(F_2\)，\(Q\) 共线，得 \(PF_1\perp PF_2\). 椭圆定义给出
\[
2a=PF_1+PF_2=(2+\sqrt2)+(2-\sqrt2)=4
\]
所以 \(a=2\). 在直角三角形 \(PF_1F_2\) 中，\(F_1F_2=2c\)，故
\[
4c^2=(2+\sqrt2)^2+(2-\sqrt2)^2=12
\]
从而 \(c=\sqrt3\)，\(b^2=a^2-c^2=1\). 因此椭圆标准方程为
\[
\frac{x^2}{4}+y^2=1
\]
\item 设 \(PF_1=PQ=n\)，\(PF_2=m\). 由椭圆定义，\(n+m=2a\). 由于 \(F_2\) 在线段 \(PQ\) 上，且 \(PQ\perp PF_1\)，在直角三角形 \(PF_1Q\) 中
\(QF_1=\sqrt{PF_1^2+PQ^2}=\sqrt2n\)，\(qquad QF_2=PQ-PF_2=n-m\).
再次利用椭圆定义，\(\sqrt2n+n-m=2a=n+m\)，所以 \(m=\frac{\sqrt2}{2}n\). 又由直角三角形 \(PF_1F_2\)，
\[
n^2+m^2=(2c)^2
\]
代入 \(m=\frac{\sqrt2}{2}n\) 得 \(n^2=\frac83c^2\). 另一方面，\(n+m=2a\) 给出
\[
n=\frac{2\sqrt2}{\sqrt2+1}a
\]
于是
\[
\frac83e^2=\frac{n^2}{a^2}=\frac{8}{(\sqrt2+1)^2}
\]
即 \(e^2=\frac{3}{3+2\sqrt2}=9-6\sqrt2=(\sqrt6-\sqrt3)^2\). 因为 \(e>0\)，故 \(e=\sqrt6-\sqrt3\).
\end{enumerate}
\end{solution}

\begin{problem}
在数列 \(\{a_{n}\}\) 中，\(a_{1} = 3\)，\(a_{n + 1}a_{n} + \lambda a_{n + 1} + \mu a_{n}^{2} = 0 \)（\(n \in \mathbf{N}_{+}\)）.

\begin{enumerate}
\item 若 \(\lambda = 0\)，\(\mu = -2\)，求数列 \(\{a_n\}\) 的通项公式；
\item 若 \(\lambda = \frac{1}{k_0}\)（\(k_0 \in \mathbf{N}_+\)，\(k_0 \geqslant 2\)），\(\mu = -1\)，证明： \(2 + \frac{1}{3k_0 + 1} < a_{k_0 + 1} < 2 + \frac{1}{2k_0 + 1}\).
\end{enumerate}
\end{problem}

\begin{solution}
\begin{enumerate}
\item 当 \(\lambda=0\)、\(\mu=-2\) 时，递推式化为
\[
a_{n+1}a_n-2a_n^2=a_n(a_{n+1}-2a_n)=0
\]
因为 \(a_1=3\ne0\)，若 \(a_n\ne0\)，则上式必有 \(a_{n+1}=2a_n\ne0\). 由数学归纳法，所有 \(a_n\ne0\)，从而对每个 \(n\) 都有 \(a_{n+1}=2a_n\). 因此 \(\{a_n\}\) 是首项为 \(3\)、公比为 \(2\) 的等比数列，故
\[
a_n=3\cdot2^{n-1}
\]
\item 记 \(K=k_0\). 当 \(\lambda=\frac1K\)、\(\mu=-1\) 时，递推式为
\[
a_{n+1}\left(a_n+\frac1K\right)=a_n^2
\]
从而
\[
a_n-a_{n+1}=\frac{a_n}{Ka_n+1}=\frac{1}{K+\frac1{a_n}}
\]
先证明在 \(1\leqslant n\leqslant K+1\) 时有 \(2<a_n\leqslant3\). 初值为 \(a_1=3\). 若 \(2<a_n\leqslant3\)，则
\[
0<a_n-a_{n+1}=\frac{1}{K+1/a_n}<\frac1K
\]
所以 \(a_{n+1}<a_n\leqslant3\). 若已知前 \(n\) 项都大于 \(2\)，逐项累加上式的严格不等式，便有
\[
a_n>3-\frac{n-1}{K}\geqslant2
\]
且在 \(n\leqslant K+1\) 时严格大于 \(2\). 因此用归纳法得到所述范围内的 \(2<a_n\leqslant3\).
因此对 \(1\leqslant n\leqslant K\)，
\(\frac13\leqslant\frac1{a_n}<\frac12\)，\(\frac{2}{2K+1}<a_n-a_{n+1}\leqslant\frac{3}{3K+1}\).
由于 \(K\geqslant2\)，且 \(a_2<a_1=3\)，故 \(n=2\) 时右侧不等式严格. 求和得到
\[
\frac{2K}{2K+1}<\sum_{n=1}^{K}(a_n-a_{n+1})=3-a_{K+1}<\frac{3K}{3K+1}
\]
整理得
\[
2+\frac{1}{3K+1}<a_{K+1}<2+\frac{1}{2K+1}
\]
代回 \(K=k_0\)，即得题中结论.
\end{enumerate}
\end{solution}
